\subsection {\href{https://leetcode.com/problems/construct-quad-tree/}{Construct Quad Tree (LeetCode 427)}}

\subsection*{Problem Description}

This problem asks us to construct a Quad-Tree representation from an $n \times n$ binary matrix.

A Quad-Tree is a tree data structure where each internal node has exactly four children representing the four quadrants: top-left, top-right, bottom-left, and bottom-right. Each node in the tree has two properties:

\begin{itemize}
\item \textbf{val}: A boolean value (True for 1s, False for 0s)
\item \textbf{isLeaf}: A boolean indicating whether the node is a leaf (True) or has four children (False)
\end{itemize}

The construction follows these rules:

\begin{itemize}

\item \textbf{If a region contains all the same values} (all 0s or all 1s):

\begin{itemize}
\item Create a leaf node with \texttt{isLeaf = True}
\item Set \texttt{val} to the grid's value (True for 1, False for 0)
\item The four children are set to null
\end{itemize}

\item \textbf{If a region contains mixed values} (both 0s and 1s):

\begin{itemize}
\item Create an internal node with \texttt{isLeaf = False}
\item Set \texttt{val} to any value (it doesn't matter since it's not a leaf)
\item Divide the current region into four equal quadrants
\item Recursively construct four children nodes for each quadrant
\end{itemize}

\end{itemize}

\subsection*{Solution Approach}

This solution uses a divide-and-conquer recursion to construct the Quad-Tree. Here is how the algorithm works step by step.

\textbf{Step 1: Check if the current region is uniform}

\begin{lstlisting}[language=C++,breaklines=true,breakatwhitespace=false,columns=fullflexible,basicstyle=\ttfamily\small]
bool isSame(vector<vector<int>>& grid, int r, int c, int size){
    int val = grid[r][c];
    for(int i = r; i < r + size; i++){
        for(int j = c; j < c + size; j++){
            if(grid[i][j] != val) return false;
        }
    }
    return true;
}
\end{lstlisting}

\begin{itemize}

\item We check whether all cells in the current square region have the same value.

\item If the region is uniform, we create a leaf node and stop recursion.

\end{itemize}

\textbf{Step 2: Recursively divide the region}

\begin{lstlisting}[language=C++,breaklines=true,breakatwhitespace=false,columns=fullflexible,basicstyle=\ttfamily\small]
Node* build(vector<vector<int>>& grid, int r, int c, int size){
    if(isSame(grid, r, c, size)){
        return new Node(grid[r][c], true);
    }

    int half = size / 2;

    Node* topLeft = build(grid, r, c, half);
    Node* topRight = build(grid, r, c + half, half);
    Node* bottomLeft = build(grid, r + half, c, half);
    Node* bottomRight = build(grid, r + half, c + half, half);

    return new Node(true, false, topLeft, topRight, bottomLeft, bottomRight);
}
\end{lstlisting}

\begin{itemize}

\item If the region contains mixed values, we divide the square into four equal quadrants.

\item We recursively construct Quad-Trees for each quadrant.

\item The current node becomes an internal node with four children.

\end{itemize}

\subsection*{Complexity Analysis}

\begin{itemize}

\item \textbf{Time Complexity}: $O(n^2 \log n)$ in the worst case because we may check each region multiple times during recursion.

\item \textbf{Space Complexity}: $O(n^2)$ for the recursion tree and Quad-Tree nodes.
\subsection*{Implementation}

\begin{lstlisting}[language=C++,breaklines=true,breakatwhitespace=false,columns=fullflexible,basicstyle=\ttfamily\small]
/*
// Definition for a QuadTree node.
class Node {
public:
    bool val;
    bool isLeaf;
    Node* topLeft;
    Node* topRight;
    Node* bottomLeft;
    Node* bottomRight;
    
    Node() {
        val = false;
        isLeaf = false;
        topLeft = NULL;
        topRight = NULL;
        bottomLeft = NULL;
        bottomRight = NULL;
    }
    
    Node(bool _val, bool _isLeaf) {
        val = _val;
        isLeaf = _isLeaf;
        topLeft = NULL;
        topRight = NULL;
        bottomLeft = NULL;
        bottomRight = NULL;
    }
    
    Node(bool _val, bool _isLeaf, Node* _topLeft, Node* _topRight,
         Node* _bottomLeft, Node* _bottomRight) {
        val = _val;
        isLeaf = _isLeaf;
        topLeft = _topLeft;
        topRight = _topRight;
        bottomLeft = _bottomLeft;
        bottomRight = _bottomRight;
    }
};
*/

class Solution {
public:

    bool isSame(vector<vector<int>>& grid, int r, int c, int size){
        int val = grid[r][c];
        for(int i = r; i < r + size; i++){
            for(int j = c; j < c + size; j++){
                if(grid[i][j] != val) return false;
            }
        }
        return true;
    }

    Node* build(vector<vector<int>>& grid, int r, int c, int size){
        if(isSame(grid, r, c, size)){
            return new Node(grid[r][c], true);
        }

        int half = size / 2;

        Node* topLeft = build(grid, r, c, half);
        Node* topRight = build(grid, r, c + half, half);
        Node* bottomLeft = build(grid, r + half, c, half);
        Node* bottomRight = build(grid, r + half, c + half, half);

        return new Node(true, false, topLeft, topRight, bottomLeft, bottomRight);
    }

    Node* construct(vector<vector<int>>& grid) {
        int n = grid.size();
        return build(grid, 0, 0, n);
    }
};
\end{lstlisting}

\end{itemize}